\hypertarget{algoritma-variational-quantum-eigensolver-vqe}{%
\section{Algoritma Variational Quantum Eigensolver (VQE)}\label{algoritma-variational-quantum-eigensolver-vqe}}

VQE adalah algoritma \emph{hybrid quantum-classical} yang memetakan masalah minimisasi energi ke dalam loop optimasi parameter \(\theta\).

\hypertarget{a.-alur-kerja-hybrid-secara-matematis}{%
\subsection{Alur Kerja Hybrid secara Matematis}\label{a.-alur-kerja-hybrid-secara-matematis}}

\begin{enumerate}
\def\labelenumi{\arabic{enumi}.}
\tightlist
\item
  \textbf{Preparasi State:} QPU menyiapkan state \(|\psi(\theta)\rangle = U(\theta)|0\rangle^{\otimes n}\).
\item
  \textbf{Dekomposisi Hamiltonian} \hyperref[appendix-e]{(lihat Appendix E)}: Hamiltonian \(\hat{H}\) adalah kombinasi linear suku-suku Pauli \(P_k \in \{I, X, Y, Z\}^{\otimes n}\): \[ \hat{H} = \sum_k c_k \hat{P}_k \]
\item
  \textbf{Pengukuran:} Nilai ekspektasi dihitung sebagai: \[ \langle \hat{H} \rangle_{\theta} = \sum_k c_k \langle \psi(\theta) | \hat{P}_k | \psi(\theta) \rangle \]
\item
  \textbf{Optimasi:} Optimizer klasik mencari \(\theta^*\) yang meminimalkan \(\langle \hat{H} \rangle_{\theta}\).
\end{enumerate}

\hypertarget{b.-optimasi-parameter-spsa-vs.-gradient-descent}{%
\subsection{Optimasi Parameter: SPSA vs.~Gradient Descent}\label{b.-optimasi-parameter-spsa-vs.-gradient-descent}}

Dalam sistem kuantum riil (noisy), estimasi gradien menggunakan \textit{Parameter Shift Rule} \hyperref[appendix-f]{(lihat Appendix F)} atau beda hingga (\textit{finite difference}) sangat rentan terhadap \emph{noise} karena memerlukan banyak pengukuran (\(2N\) pengukuran per iterasi).

\hypertarget{masalah-gradient-descent-gd}{%
\subsubsection{Masalah Gradient Descent (GD)}\label{masalah-gradient-descent-gd}}

GD menghitung turunan parsial terhadap setiap parameter \(\theta_i\) secara bergantian. Pada lanskap energi yang ``kasar'' akibat \emph{noise} statistik (\emph{shot noise}), GD cenderung terjebak pada \emph{local minima} atau gagal konvergen karena estimasi gradien yang bias.

\hypertarget{mekanisme-spsa-simultaneous-perturbation-stochastic-approximation}{%
\subsubsection{Mekanisme SPSA (Simultaneous Perturbation Stochastic Approximation)}\label{mekanisme-spsa-simultaneous-perturbation-stochastic-approximation}}

SPSA adalah algoritma optimasi stokastik yang jauh lebih efisien untuk VQE. Alih-alih mengukur gradien satu per satu, SPSA melakukan perturbasi pada \textbf{seluruh parameter sekaligus} menggunakan vektor acak \(\Delta\):

\begin{enumerate}
\def\labelenumi{\arabic{enumi}.}
\tightlist
\item
  \textbf{Perturbasi Acak:} Pilih vektor \(\Delta_k\) di mana setiap elemen bernilai \(\pm 1\) secara acak (distribusi Bernoulli).
\item
  \textbf{Estimasi Gradien Stokastik:} Hitung gradien pendekatan \(g(\theta_k)\) hanya dengan \textbf{dua pengukuran} fungsi biaya: \[ \hat{g}(\theta_k) = \frac{E(\theta_k + c_k \Delta_k) - E(\theta_k - c_k \Delta_k)}{2 c_k \Delta_k} \] di mana \(c_k\) adalah besarnya perturbasi yang mengecil seiring iterasi.
\item
  \textbf{Update Parameter:} \[ \theta_{k+1} = \theta_k - a_k \hat{g}(\theta_k) \]
\end{enumerate}

\hypertarget{keunggulan-spsa-untuk-markowitz-vqe}{%
\subsubsection{Keunggulan SPSA untuk Markowitz-VQE}\label{keunggulan-spsa-untuk-markowitz-vqe}}

\begin{itemize}
\tightlist
\item
  \textbf{Efisiensi Skalabilitas:} Hanya memerlukan 2 \emph{evaluasi fungsi} per iterasi, terlepas dari jumlah aset \(N\). GD memerlukan \(2N\) evaluasi.
\item
  \textbf{Robustness terhadap Noise:} SPSA memiliki sifat ``averaging'' alami yang membuatnya lebih tahan terhadap fluktuasi statistik hasil pengukuran QPU dibandingkan \emph{optimizer} berbasis gradien analitik.
\item
  \textbf{Konvergensi Global:} Karena sifat stokastiknya, SPSA memiliki peluang lebih besar untuk melompati \emph{local minima} kecil pada lanskap energi portfolio.
\end{itemize}


